\chapter{Glosarium}
\label{app:glosarium}

Glosarium ini memuat istilah yang dipakai sepanjang buku beserta
padanan bahasa Inggrisnya dan penjelasan singkat. Istilah disusun
menurut abjad Indonesia.

\begingroup
\footnotesize
\begin{longtable}{@{}p{3.4cm}p{3.4cm}p{6.6cm}@{}}
\toprule
\textbf{Istilah} & \textbf{Padanan} & \textbf{Penjelasan} \\
\midrule
\endfirsthead
\toprule
\textbf{Istilah} & \textbf{Padanan} & \textbf{Penjelasan} \\
\midrule
\endhead
\bottomrule
\endfoot

Alamat efektif & \eng{effective address} & Alamat yang dihasilkan
dari perhitungan base + index + displacement pada mode pengalamatan
tertentu. \\
\addlinespace
Alamat fisik & \eng{physical address} & Alamat 20-bit yang benar-benar
dikirim ke bus alamat, hasil perhitungan segmen dan offset. \\
\addlinespace
Assembler & \eng{assembler} & Program yang menerjemahkan kode sumber
assembly menjadi kode mesin dan berkas objek. \\
\addlinespace
Bilangan bertanda & \eng{signed number} & Bilangan yang dapat bernilai
negatif, disandikan dengan \eng{two's complement} pada x86. \\
\addlinespace
Digit biner & \eng{bit} & Satuan informasi terkecil, bernilai 0 atau 1. \\
\addlinespace
Instruksi mesin & \eng{machine instruction} & Instruksi dalam bentuk
deretan bit yang langsung dikenali prosesor. \\
\addlinespace
Interupsi & \eng{interrupt} & Mekanisme penghentian sementara alur
program untuk menjalankan layanan sistem. \\
\addlinespace
Kode mesin & \eng{machine code} & Program dalam bentuk deretan bit. \\
\addlinespace
Kode sumber & \eng{source code} & Program dalam bentuk teks yang
ditulis programmer. \\
\addlinespace
Lompatan bersyarat & \eng{conditional jump} & Instruksi lompatan yang
hanya dijalankan bila kondisi flag tertentu terpenuhi. \\
\addlinespace
Memori konvensional & \eng{conventional memory} & 640 KB pertama
ruang memori yang dapat dipakai program DOS. \\
\addlinespace
Mode pengalamatan & \eng{addressing mode} & Cara penulisan operand
yang menentukan bagaimana alamat operand dihitung. \\
\addlinespace
Mnemonic & \eng{mnemonic} & Nama pendek yang mewakili sebuah opcode
dalam bahasa assembly, misalnya \asm{MOV}. \\
\addlinespace
Nibble & \eng{nibble} & Kelompok empat bit, setara dengan satu digit
heksadesimal. \\
\addlinespace
Opcode & \eng{opcode} & Bagian instruksi yang menyatakan operasi yang
harus dilakukan. \\
\addlinespace
Operand & \eng{operand} & Bagian instruksi yang menyatakan data yang
dikenai operasi. \\
\addlinespace
Overflow & \eng{overflow} & Kelebihan hasil pada operasi bilangan
bertanda; ditandai flag \flag{OF}. \\
\addlinespace
Paragraf & \eng{paragraph} & Satuan 16 byte, yaitu kelipatan terkecil
yang dapat ditunjuk sebuah segmen. \\
\addlinespace
Penyangga & \eng{buffer} & Area memori sementara untuk menampung data
masukan atau keluaran. \\
\addlinespace
Perakitan & \eng{assembling} & Proses menerjemahkan kode sumber
assembly menjadi berkas objek. \\
\addlinespace
Pergeseran aritmetika & \eng{arithmetic shift} & Pergeseran yang
mempertahankan bit tanda pada bilangan bertanda. \\
\addlinespace
Pergeseran logis & \eng{logical shift} & Pergeseran yang mengisi bit
kosong dengan nol. \\
\addlinespace
Pipelining & \eng{pipelining} & Teknik menjalankan tahapan instruksi
berbeda secara tumpang tindih untuk menaikkan kinerja. \\
\addlinespace
Port I/O & \eng{port I/O} & Akses perangkat melalui ruang alamat
terpisah memakai instruksi \asm{IN} dan \asm{OUT}. \\
\addlinespace
Register & \eng{register} & Penyimpan kecil berkecepatan tinggi di
dalam prosesor. \\
\addlinespace
Register flag & \eng{flag register} & Register yang menampung bit-bit
status hasil operasi sebelumnya. \\
\addlinespace
Rekursi & \eng{recursion} & Prosedur yang memanggil dirinya sendiri. \\
\addlinespace
Rotasi & \eng{rotate} & Pergeseran yang memindahkan bit yang keluar
kembali ke ujung yang lain. \\
\addlinespace
Segmen & \eng{segment} & Blok memori berukuran maksimum 64 KB yang
ditunjuk sebuah register segmen. \\
\addlinespace
Siklus instruksi & \eng{instruction cycle} & Rangkaian
\eng{fetch}, \eng{decode}, dan \eng{execute} untuk satu instruksi. \\
\addlinespace
Siklus mesin & \eng{machine cycle} & Satu kali pengaksesan bus oleh
prosesor. \\
\addlinespace
Simbol & \eng{symbol} & Nama yang ditetapkan assembler untuk sebuah
label, variabel, atau konstanta. \\
\addlinespace
Stack & \eng{stack} & Area memori yang dikelola secara
\eng{LIFO} untuk menyimpan alamat kembali dan data sementara. \\
\addlinespace
Stack frame & \eng{stack frame} & Susunan data di stack selama satu
pemanggilan prosedur, biasanya ditunjuk register \reg{BP}. \\
\addlinespace
Tabel vektor interupsi & \eng{interrupt vector table} & Tabel pada
alamat 0000h--03FFh yang memuat alamat setiap penangan interupsi. \\
\addlinespace
Tabel simbol & \eng{symbol table} & Daftar nama simbol beserta
alamatnya, dihasilkan assembler. \\
\addlinespace
Two's complement & \eng{two's complement} & Penyandian bilangan
bertanda dengan membalik bit lalu menambahkan satu. \\
\addlinespace
\end{longtable}
\endgroup
